% Templat TA TIF UIN Suska Riau -- (c) 2026 Benny Sukma Negara
\chapter{Metodologi Penelitian}
\label{bab:metodologi}

%==========================================================
\section{Tahapan Penelitian}
\label{sec:tahapan}

Bagian metodologi berisi penjelasan mengenai hipotesa, variabel-variabel, model, rancangan penelitian, teknik pengumpulan data, teknik analisis data, cara penafsiran dan penyimpulan hasil penelitian yang digunakan \parencite{sugiyono2019metode}. Tahapan penelitian ditunjukkan pada Gambar~\ref{fig:tahapan}.

\begin{figure}[htbp]
\centering
\gambar[0.7\textwidth]{tahapan_penelitian.png}
\caption{Tahapan penelitian}
\label{fig:tahapan}
\end{figure}

%==========================================================
\section{Pengumpulan Data}
\label{sec:data}

\isi{Sumber data (primer/sekunder), populasi dan sampel, instrumen, periode pengambilan, dan pra-pemrosesan data.}

%==========================================================
\section{Analisis dan Perancangan}
\label{sec:rancangan}

\isi{Analisis kebutuhan dan perancangan sistem/model (mis. \textit{Waterfall}, \textit{Prototyping}, CRISP-DM), diagram UML/ERD/arsitektur, atau rancangan eksperimen.}

%==========================================================
\section{Pengujian dan Evaluasi}
\label{sec:pengujian}

\isi{Metode dan metrik evaluasi (mis. \textit{black-box testing}, SUS, akurasi/presisi/\textit{recall}/F1) serta kriteria keberhasilan.}

\iflaporan\else
%==========================================================
\section{Jadwal Penelitian}
\label{sec:jadwal}

Jadwal pelaksanaan penelitian disajikan pada Tabel~\ref{tab:jadwal}.

\begin{table}[htbp]
\caption{Jadwal penelitian}
\label{tab:jadwal}
\begin{tabular}{|l|*{6}{c|}}
\hline
\multirow{2}{*}{Kegiatan} & \multicolumn{6}{c|}{Bulan ke-}\\
\cline{2-7}
 & 1 & 2 & 3 & 4 & 5 & 6\\
\hline
Studi literatur            & $\checkmark$ & $\checkmark$ &  &  &  & \\
\hline
Pengumpulan data           &  & $\checkmark$ & $\checkmark$ &  &  & \\
\hline
Analisis dan perancangan   &  &  & $\checkmark$ & $\checkmark$ &  & \\
\hline
Implementasi               &  &  &  & $\checkmark$ & $\checkmark$ & \\
\hline
Pengujian dan evaluasi     &  &  &  &  & $\checkmark$ & $\checkmark$\\
\hline
Penyusunan laporan         & $\checkmark$ & $\checkmark$ & $\checkmark$ & $\checkmark$ & $\checkmark$ & $\checkmark$\\
\hline
\end{tabular}
\end{table}
\fi
